% THIS IS A LATEX TEMPLATE FILE FOR PAPERS INCLUDED IN THE
% *Anthology of Computers and the Humanities*.
\documentclass[final]{anthology-ch}

% LOAD LaTeX PACKAGES
\usepackage{booktabs}
\usepackage{graphicx}

% TITLE OF THE SUBMISSION
\title{From Working Group to Campus Resource: Strengthening Digital Humanities through the Library}

% AUTHOR AND AFFILIATION INFORMATION
\author[1]{Christin Lampkowski}[
  orcid=0009-0005-6929-1891
]

\author[1]{Savannah Lake}[
  orcid=0000-0001-9998-5185
]

\author[1]{Jenna Rinalducci}[
  orcid=0000-0002-7607-5653
]

\affiliation{1}{University of North Carolina at Charlotte, Charlotte, NC, USA}

% KEYWORDS
\keywords[eng]{digital humanities, academic libraries, library working group, digital scholarship, communities of practice, institutional support}

% METADATA FOR THE PUBLICATION
\pubyear{2025}
\pubvolume{5}
\pagestart{63}
\pageend{69}
\paperorder{8}
\conferencename{From Manuscript to Megabyte: Building Sustainable and Open Futures in the Digital Humanities}
\conferenceeditors{James B. Harr III, Emma Stanley, Jenna Rinalducci, and Donna Kain}
\doi{10.63744/NN0FNCnEfR92}
\customhead{}

\addbibresource{bibliography.bib}

%%%%%%%%%%%%%%%%%%%%%%%%%%%%%%%%%%%%%%%%%%%%%%%%%%%%%%%%%%%%%%%%%%%%%%%%%%%
% HERE IS THE START OF THE TEXT
\begin{document}

\maketitle

\begin{abstract}
Digital humanities activity at the University of North Carolina at Charlotte has expanded significantly in recent years, yet this growth has occurred without a centralized hub or dedicated staffing. To address some of these gaps, J. Murrey Atkins Library created a Digital Humanities Working Group to coordinate support, develop programming, and build community across campus. This article outlines the evolution of the working group, examines findings from a recent environmental scan, and describes current initiatives and future directions. The library's position as a connector across disciplines provides a strategic foundation for strengthening digital humanities research, teaching, and collaboration at UNC Charlotte.
\end{abstract}

\section{Introduction} \label{sec:intro}

Much of the existing scholarship focuses on well-resourced institutions or formal digital humanities centers, leaving fewer models for campuses where digital humanities activity is active but decentralized. However, digital humanities research and instruction frequently develop through localized initiatives rather than through centralized programs, as seen at the University of North Carolina at Charlotte (UNC Charlotte). Within this landscape, academic libraries have emerged as important contributors to digital humanities work. Libraries bring expertise in research support, digital scholarship, metadata, preservation, and instructional collaboration, and they often serve as interdisciplinary spaces where faculty and students from different departments can connect.

Almost ten years ago, an informal gathering of librarians with a shared interest in digital humanities gradually developed into a coordinated, interdepartmental effort that supports work across campus. What started as a discussion group gradually became the Digital Humanities Working Group, a collaborative committee representing departments across the library, including Public Services, Digital Scholarship and Innovation, and Special Collections and University Archives. These efforts led to a formalized library group focused on identifying existing digital humanities activity across campus, improving communication channels, and developing programming that could support faculty and students in a variety of disciplines. The shift reflected both the growing interest in digital humanities on campus and the library's role as a connector across disciplines.

\section{Literature Review} \label{sec:litreview}

Numerous case studies and articles have explored how librarians are well suited to be instructors, collaborators, and facilitators of digital humanities work. This is in part because librarians are already experts in practices that are naturally a part of digital humanities work, including digitization, metadata, copyright, and preservation, leading to successful collaborative teaching between librarians and faculty.\footnote{See, for example, multiple case studies in \cite{hartsellgundy2015digital}, including \cite[69--82]{borovsky2015digital} and \cite[151--176]{rosenblum2015collaboration}.} In addition to their expertise, librarians also occupy a unique position at academic institutions by engaging in what Bell and Kennan describe as a \enquote{liminal space,} complicating recognition as staff, academic, and scholar, but providing \enquote{greater opportunity for mediation between academic faculty, students, and non-academic communities represented in DH collections and library spaces}~\cite[pp.~166--167]{bell2021partnering}. By not being tied to one department, librarians are well equipped to connect digital humanities researchers across campus, who may come from traditional humanities disciplines as well as the sciences and social sciences.

Library support and participation in digital humanities work is evident in digital humanities research output as well. A quantitative analysis of authorship in digital humanities journals found that digital humanities scholarship is highly collaborative, with nearly 45 percent of articles written by multiple authors, averaging to 3.2 authors per coauthored article. As compared with eight fields across the social sciences and sciences (to include economics, biomedical research, mathematics, physics, computer science, information systems, and sociology), these co-authorship metrics rank third and second, respectively~\cite[p.~682]{wong2016digital}. The professions of the authors in the study are largely scholars (31 percent), with digital humanities center staff and library staff co-authors representing 10 percent and 9 percent of co-authors, respectively~\cite[p.~679]{wong2016digital}.

While it is clear how librarians and libraries would be natural contributors to digital humanities work at universities, the library's role in digital humanities work can look different from university to university. Institutions have varying organizational structures, financial resources, priorities, and cultures of practice that shape how digital humanities support and research is organized on campus~\cite[pp.~43--45]{posner2013half}. Important, too, is the varying levels of involvement from other campus collaborators, such as IT and various departments~\cite[p.~11]{anne2017building}. The approach some institutions have taken is the creation of a digital humanities center for their campus. These centers are often established in collaboration with libraries, even if they are not housed within the library~\cite{white2016preface}. In other cases, libraries provide support such as software and technology access, workshops or classroom instruction, and collaborative support without having a physical space or center designated for digital humanities~\cite{habing2025digital}.

Institutions facing funding constraints and a lack of staffing have explored a number of approaches to digital humanities support. Often this takes the form of \enquote{ad hoc} support, with \enquote{the success of DH efforts in the library often depend[ing] on the energy, creativity, and goodwill of a few overextended library and IT professionals and the services they can throw together}~\cite{bellc2015practice}. At Salem State University, the library partnered with an English faculty member and utilized an internal university grant to develop digital humanities support, including a project to house digital scholarship and an undergraduate research program---both of which served as \enquote{proofs of concept that make the case for institutional investment in them and in digital humanities center}~\cite[p.~348]{risam2017building}.

Digital humanities centers at colleges and universities take many different forms, but several clear trends define their current landscape. Institutions increasingly treat digital humanities as an interdisciplinary space that supports research, teaching, and collaborative project work, yet the level of investment and infrastructure varies widely. Some campuses, such as San Diego State University, maintain staffed physical centers within the library that include studios and flexible gathering spaces that host workshops, instruction, and project development~\cite{lach2024interview}. However, many centers face ongoing staffing challenges and post-pandemic retention issues that affect long-term stability. Others, like the University of Nebraska--Lincoln, offer a more expansive model that includes undergraduate and graduate credentials, a steady slate of events, and a substantial portfolio of digital scholarship projects, positioning them as aspirational programs within the field~\cite{cdrh2026anniversary}. Even institutions without formal centers contribute through distributed support models, including online guides like LibGuides, toolkits, and cross-unit collaborations that help sustain digital humanities activity. Taken together, these examples show a field that is growing and diversifying, with strong models in place but ongoing needs for sustainable staffing, institutional commitment, and adaptable support structures.

\section{Current Practices} \label{sec:current}

Digital humanities research at UNC Charlotte is largely decentralized, with researchers coming from departments across campus and no central institute or center supporting digital humanities work. In the absence of a central unit to coordinate activity, digital humanities work often develops independently within departments, and faculty rely on varied interpretations of what digital humanities entails. This distributed structure also limits funding for technical support and staffing, and it places stakeholders across multiple colleges, departments, and units, including the library. While many faculty on campus have and continue to do digital humanities work, in 2025 the university hired the first faculty position to specify digital humanities work in the position title and description, a joint faculty position between the Department of English and School of Data Science. Aside from this position, though, the campus and its researchers lack formalized digital humanities staff and support.

Working without a centralized digital humanities hub requires the library to approach this work with limited staffing and distributed responsibilities. These constraints shape how we consider the support we can realistically provide. Although these conditions present obstacles, they also position the library to help connect practitioners, support emerging projects, and strengthen communication across campus through networking, skills development, and the cultivation of communities of practice. As the working group began to formalize its efforts, several guiding questions emerged that continue to inform our planning. We asked how we could support digital humanities research without dedicated staff or a center, how we might build relationships with practitioners across campus, how to make connections when definitions of digital humanities vary, and how to share ongoing work in ways that strengthen visibility and impact. These questions help us identify where the library can contribute most effectively and where collaboration is essential. In response to these considerations, the working group identified several focus areas that align with the library's strengths and its role as a connector on campus. These include collaboration, facilitation, strategic partnerships, communication, and advocacy. Together, these areas provide a framework for supporting digital humanities in a distributed environment and guide the initiatives described in the following sections.

Given these focus areas, much of the working group's current work revolves around creating opportunities for connection and developing communities of practice. The library has established two platforms for communication. The first is a campuswide discussion group that anyone interested in digital humanities can join. This Google Group is intended for faculty and students to ask questions about the digital humanities and their projects, seek collaborators, and share resources and opportunities, such as calls for proposals and grants. The second mode of communication is an email address specific to the Digital Humanities Library Working Group, where faculty and students can reach out to librarians for digital humanities support.

The events that the Digital Humanities Library Working Group hosts similarly have evolved to focus on building connections and communities of practice. In the group's early years, from 2019 to 2022, the focus for events was largely on hosting workshops, webinars, and conferences. The limitation of this approach was that these events required substantial planning, and given the group's limited resources, the result was that these events were sporadic. Since the fall of 2024, the group has transitioned to a new strategy for events, where once a semester we host a digital humanities campus meet-up. This is an opportunity for digital humanities practitioners across campus to informally meet up. Generally, we will have a theme or discussion topic for the meet-up, while also letting conversations naturally go where people are interested. Since these events require less planning, we are able to host them more regularly, which we hope will help people develop community and connections.

One resource that the library has been able to leverage in service of digital humanities support is the library's fellowship program. Atkins Library is fortunate to host summer fellows, where graduate students in library science programs complete a ten-week fellowship. Members of the group have worked with these fellows, to create both library-based digital humanities projects and new digital humanities support for the campus at large. Library-based projects include a digital exhibit on artist books housed in special collections~\cite{atkinsartists} as well as a digital timeline of LGBTQIA+ history in Charlotte, showcasing materials and gaps in the library's collections. Campus-facing projects include the creation of the library's 3D modelling and photogrammetry service as well as a 2024 report on how the library can better support digital humanities drawn from surveys and reviews of resources on campus and at peer institutions~\cite{anderson2024facilitating}. This report importantly helped us evaluate our current services against faculty priorities, orienting our future work as a group.

\section{Future Directions} \label{sec:future}

As the Digital Humanities Working Group continues to evolve, the 2024 report on assessing digital humanities needs on campus against the support offered at peer and aspirational institutions has been instrumental in informing our future directions. Completed largely through a ten-week fellowship, the report draws from findings from interviews with twenty-one people, including eight teaching faculty, eight library staff, and five people from outside of UNC Charlotte, to include faculty affiliated with other University of North Carolina System schools, schools with established digital humanities centers, and the National Endowment for the Humanities~\cite[p.~3]{anderson2024facilitating}. Feedback from faculty on campus shows that the most requested resource is access to tools, followed closely by education, space, communication, and staffing (Figure~\ref{fig:resources})~\cite[p.~9]{anderson2024facilitating}.

\begin{figure}[t!]
  \centering
  \includegraphics[width=0.9\linewidth]{images/figure1.png}
  \caption{Desired faculty resources for DH support, from the 2024 DH report.}
  \label{fig:resources}
\end{figure}

As much as the library would like to support digital humanities research in these ways, having dedicated space and staffing requires substantial financial investments. Increased tools and education also require a certain level of staffing and financial support. In reviewing how digital humanities centers and libraries offer such services, we found a number of staffing and funding models. One model is through college or departmental sponsorships to digital humanities centers, wherein colleges or departments offer either funding or tuition remission for student worker positions. The Institute for Digital Research in the Humanities at the University of Kansas, for example, is \enquote{a partnership between KU Libraries, The Hall Center for the Humanities, and The College of Liberal Arts and Sciences}~\cite{kansashistory}. Another financial model leverages research grant funding, with digital humanities projects hosted at digital humanities centers utilizing external grants~\cite{gayman2019digital}. Support from the university at large, or from private donors, has also been used to build out digital humanities centers. At George Mason University, the Roy Rosenzweig Center for History and New Media, for example, relies on a number of revenue streams, including private donors, funding from the National Endowment for the Humanities, and an affiliation with the Department of History and Art History~\cite{rrchnmstory}.

Currently, UNC Charlotte and Atkins Library are not well positioned to implement any of the above funding models for establishing a digital humanities center with staff. However, the 2024 DH report did identify some faculty needs for digital humanities support that the library is investigating further. A common sentiment expressed by faculty in the report on digital humanities needs on campus revolved around increased support for digital humanities education, communication, and networking~\cite[p.~2]{anderson2024investigating}. Given this, the group is exploring how we can create a digital hub for learning and networking. Currently, the group maintains a resource guide on digital humanities~\cite{atkinsdh}, complete with information on digital humanities research methods and tools, relevant literature and news, on-campus resources and projects, and professional networks and opportunities. The group is exploring how to transform this resource guide into a digital hub that is more current, supportive, and dynamic, with stakeholders across the campus having a role in its management. While having a physical center or space for digital humanities may not be financially tenable, a digital hub for this work may be able to provide a sense of community and education. Currently, we are developing a website for this purpose. We anticipate that this robust, flexible platform---which is less static than the platform for the current research guide---will allow us to better showcase resources, tools, projects, and community, providing a collective vision for the digital humanities at UNC Charlotte.

As changes at the national level create an even more competitive environment for accessing funding for digital humanities work, the need for campus-based funding to support work in this area increased. The library remains a logical centralized location to make campus-funded investment in software and educational opportunities that multiple disciplines can access. The library's digital humanities working group is well positioned to identify software, workshops, and trainings that are of interest across departments and advocate for campus financial support to host these. Software available in the library and events hosted by us are available to all faculty, staff, and students on campus, which also increases real-world networking opportunities across departments, sparking potential for cross-departmental collaboration on future digital humanities work. New library leadership has also presented renewed energy to advocate for the work librarians do related to digital humanities, through information sharing and funding conversations.

Student engagement with digital humanities is another area where the group sees potential for expansion. We hope to increase student participation in future workshops, opening up these opportunities beyond graduate students to include undergraduate students as well. We see this as an opportunity for undergraduate students to learn more about digital humanities, piquing their interest through skill building with specific tools in an environment that exposes them to the potential uses in digital humanities projects and interactions with upper-level students or faculty with whom they can study or collaborate. The digital hub will serve as a location to share these and other education opportunities, such as digital humanities related courses, with students. There is also potential to highlight student projects in the hub to increase visibility of student work, not just student support of faculty projects. The library is also considering student employment opportunities related to digital humanities, starting with the potential for a student position to work developing the digital hub to create relevant content for building student interest.

\section{Conclusion} \label{sec:conclusion}

As practitioners at many institutions have highlighted, there are numerous ways to support digital humanities work on university campuses. While digital humanities centers or dedicated digital humanities staff are an approach taken by many institutions, there is no one-size-fits-all approach that works everywhere. At UNC Charlotte, library faculty continue to see the library as the logical nexus for digital humanities work on our campus. The interdisciplinary nature of the library means we are the natural place where faculty and students from different departments and interest areas can cross paths, accessing educational and networking opportunities, technology, and expertise to build a more cohesive community of practice around digital humanities on our campus. By continuing to build relationships with faculty outside the library that work and teach in the digital humanities, we hope to grow the in-person and online community connections and expand educational opportunities. Though the approaches taken by librarians continue to evolve, our commitment remains unwavering.

\printbibliography

\end{document}
